\section{十九世纪冲突其余账本事件锚点}
本节补齐1851至1874年余下账本条目，保持军政叙事、外交文本、地方档案与战后社会影响的证据界限。
\subsection{1850—1859}
\eventanchor{EQ-1853-NANJING-TAKING-DECISION}{围绕「太平军攻占南京并改名天…}咸丰三年（1853年），围绕「太平军攻占南京并改名天京，建立长期政权中心」的命令、调度、照会或呈报形成独立的中央—地方行动文书链。核心取证：《清文宗实录》咸丰三年相关条；军机处档、战事方略及地方奏报。该项锚定的是命令或决策环节，命令抵达和结果仍须另外核验。来源保留：奏折、上谕、部议、军机档或条约可证明行政动作和机构立场，不自动证明所有地区、群体或对手按其意图行动。
\eventanchor{EQ-1853-NANJING-TAKING-AFTERMATH}{「太平军攻占南京并改名天京，…}咸丰三年（1853年），「太平军攻占南京并改名天京，建立长期政权中心」引发的善后、执行或地方回应构成可由不同材料追踪的后续节点。核心取证：《清文宗实录》咸丰三年相关条；军机处档、战事方略及地方奏报；相关地区的档案、志书、契约、报刊或对方材料。该项锚定的是后续追踪问题，影响范围须按地区与群体分别判断。来源保留：战后、改革后或条约后的记忆、地方记录和官方总结常具有事后选择性；后续影响须按地点、群体和时间分层，不作整体化推断。
\eventanchor{EQ-1853-NORTHERN-EXPEDITION-PRELUDE}{「太平军发动北伐，进逼华北但…}咸丰三年（1853年），「太平军发动北伐，进逼华北但补给与指挥困难加重」之前的条件、信息和争议已通过中央与地方材料显现，构成该事件可追踪的前史节点。核心取证：《清文宗实录》咸丰三年相关条；军机处档、战事方略及地方奏报；同时期地方档案、地方志、日记或对方材料应按具体地区补检。该项锚定的是前史条件与信息背景，不把条件直接写成唯一因果。来源保留：官修编年和地方材料的记录密度与立场并不相同；本行不将条件、传闻、呈报或后见解释直接等同为确定因果。
\eventanchor{EQ-1853-NORTHERN-EXPEDITION-DECISION}{围绕「太平军发动北伐，进逼华…}咸丰三年（1853年），围绕「太平军发动北伐，进逼华北但补给与指挥困难加重」的命令、调度、照会或呈报形成独立的中央—地方行动文书链。核心取证：《清文宗实录》咸丰三年相关条；军机处档、战事方略及地方奏报。该项锚定的是命令或决策环节，命令抵达和结果仍须另外核验。来源保留：奏折、上谕、部议、军机档或条约可证明行政动作和机构立场，不自动证明所有地区、群体或对手按其意图行动。
\eventanchor{EQ-1853-NORTHERN-EXPEDITION-AFTERMATH}{「太平军发动北伐，进逼华北但…}咸丰三年（1853年），「太平军发动北伐，进逼华北但补给与指挥困难加重」引发的善后、执行或地方回应构成可由不同材料追踪的后续节点。核心取证：《清文宗实录》咸丰三年相关条；军机处档、战事方略及地方奏报；相关地区的档案、志书、契约、报刊或对方材料。该项锚定的是后续追踪问题，影响范围须按地区与群体分别判断。来源保留：战后、改革后或条约后的记忆、地方记录和官方总结常具有事后选择性；后续影响须按地点、群体和时间分层，不作整体化推断。
\eventanchor{EQ-1853-WESTERN-EXPEDITION-PRELUDE}{「太平军西征争夺长江中上游，…}咸丰三年（1853年），「太平军西征争夺长江中上游，战事牵动安徽、江西、湖北」之前的条件、信息和争议已通过中央与地方材料显现，构成该事件可追踪的前史节点。核心取证：《清文宗实录》咸丰三年相关条；军机处档、战事方略及地方奏报；同时期地方档案、地方志、日记或对方材料应按具体地区补检。该项锚定的是前史条件与信息背景，不把条件直接写成唯一因果。来源保留：官修编年和地方材料的记录密度与立场并不相同；本行不将条件、传闻、呈报或后见解释直接等同为确定因果。
\eventanchor{EQ-1853-WESTERN-EXPEDITION-DECISION}{围绕「太平军西征争夺长江中上…}咸丰三年（1853年），围绕「太平军西征争夺长江中上游，战事牵动安徽、江西、湖北」的命令、调度、照会或呈报形成独立的中央—地方行动文书链。核心取证：《清文宗实录》咸丰三年相关条；军机处档、战事方略及地方奏报。该项锚定的是命令或决策环节，命令抵达和结果仍须另外核验。来源保留：奏折、上谕、部议、军机档或条约可证明行政动作和机构立场，不自动证明所有地区、群体或对手按其意图行动。
\eventanchor{EQ-1853-WESTERN-EXPEDITION-AFTERMATH}{「太平军西征争夺长江中上游，…}咸丰三年（1853年），「太平军西征争夺长江中上游，战事牵动安徽、江西、湖北」引发的善后、执行或地方回应构成可由不同材料追踪的后续节点。核心取证：《清文宗实录》咸丰三年相关条；军机处档、战事方略及地方奏报；相关地区的档案、志书、契约、报刊或对方材料。该项锚定的是后续追踪问题，影响范围须按地区与群体分别判断。来源保留：战后、改革后或条约后的记忆、地方记录和官方总结常具有事后选择性；后续影响须按地点、群体和时间分层，不作整体化推断。
\eventanchor{EQ-1854-TAIPING-NORTHERN-DEFEAT-PRELUDE}{「太平军北伐受挫，清军逐步恢…}咸丰四年（1854年），「太平军北伐受挫，清军逐步恢复华北防线」之前的条件、信息和争议已通过中央与地方材料显现，构成该事件可追踪的前史节点。核心取证：《清文宗实录》咸丰四年相关条；军机处档、战事方略及地方奏报；同时期地方档案、地方志、日记或对方材料应按具体地区补检。该项锚定的是前史条件与信息背景，不把条件直接写成唯一因果。来源保留：官修编年和地方材料的记录密度与立场并不相同；本行不将条件、传闻、呈报或后见解释直接等同为确定因果。
\eventanchor{EQ-1854-TAIPING-NORTHERN-DEFEAT-DECISION}{围绕「太平军北伐受挫，清军逐…}咸丰四年（1854年），围绕「太平军北伐受挫，清军逐步恢复华北防线」的命令、调度、照会或呈报形成独立的中央—地方行动文书链。核心取证：《清文宗实录》咸丰四年相关条；军机处档、战事方略及地方奏报。该项锚定的是命令或决策环节，命令抵达和结果仍须另外核验。来源保留：奏折、上谕、部议、军机档或条约可证明行政动作和机构立场，不自动证明所有地区、群体或对手按其意图行动。
\eventanchor{EQ-1854-TAIPING-NORTHERN-DEFEAT-AFTERMATH}{「太平军北伐受挫，清军逐步恢…}咸丰四年（1854年），「太平军北伐受挫，清军逐步恢复华北防线」引发的善后、执行或地方回应构成可由不同材料追踪的后续节点。核心取证：《清文宗实录》咸丰四年相关条；军机处档、战事方略及地方奏报；相关地区的档案、志书、契约、报刊或对方材料。该项锚定的是后续追踪问题，影响范围须按地区与群体分别判断。来源保留：战后、改革后或条约后的记忆、地方记录和官方总结常具有事后选择性；后续影响须按地点、群体和时间分层，不作整体化推断。
\eventanchor{EQ-1854-XIANG-ARMY-ORGANIZATION-PRELUDE}{「湘军编制、饷源和将领网络逐…}咸丰四年（1854年），「湘军编制、饷源和将领网络逐步成形」之前的条件、信息和争议已通过中央与地方材料显现，构成该事件可追踪的前史节点。核心取证：《清文宗实录》咸丰四年相关条；宫中朱批奏折、内阁大库题本与《清史稿》相关志传；同时期地方档案、地方志、日记或对方材料应按具体地区补检。该项锚定的是前史条件与信息背景，不把条件直接写成唯一因果。来源保留：官修编年和地方材料的记录密度与立场并不相同；本行不将条件、传闻、呈报或后见解释直接等同为确定因果。
\eventanchor{EQ-1854-XIANG-ARMY-ORGANIZATION-DECISION}{围绕「湘军编制、饷源和将领网…}咸丰四年（1854年），围绕「湘军编制、饷源和将领网络逐步成形」的命令、调度、照会或呈报形成独立的中央—地方行动文书链。核心取证：《清文宗实录》咸丰四年相关条；宫中朱批奏折、内阁大库题本与《清史稿》相关志传。该项锚定的是命令或决策环节，命令抵达和结果仍须另外核验。来源保留：奏折、上谕、部议、军机档或条约可证明行政动作和机构立场，不自动证明所有地区、群体或对手按其意图行动。
\eventanchor{EQ-1854-XIANG-ARMY-ORGANIZATION-AFTERMATH}{「湘军编制、饷源和将领网络逐…}咸丰四年（1854年），「湘军编制、饷源和将领网络逐步成形」引发的善后、执行或地方回应构成可由不同材料追踪的后续节点。核心取证：《清文宗实录》咸丰四年相关条；宫中朱批奏折、内阁大库题本与《清史稿》相关志传；相关地区的档案、志书、契约、报刊或对方材料。该项锚定的是后续追踪问题，影响范围须按地区与群体分别判断。来源保留：战后、改革后或条约后的记忆、地方记录和官方总结常具有事后选择性；后续影响须按地点、群体和时间分层，不作整体化推断。
\eventanchor{EQ-1855-YELLOW-RIVER-COURSE-SHIFT-PRELUDE}{「黄河在铜瓦厢决口改道北流，…}咸丰五年（1855年），「黄河在铜瓦厢决口改道北流，重塑华北水系与社会经济」之前的条件、信息和争议已通过中央与地方材料显现，构成该事件可追踪的前史节点。核心取证：《清文宗实录》咸丰五年相关条；地方志、督抚奏折及地方档案；同时期地方档案、地方志、日记或对方材料应按具体地区补检。该项锚定的是前史条件与信息背景，不把条件直接写成唯一因果。来源保留：官修编年和地方材料的记录密度与立场并不相同；本行不将条件、传闻、呈报或后见解释直接等同为确定因果。
\eventanchor{EQ-1855-YELLOW-RIVER-COURSE-SHIFT-DECISION}{围绕「黄河在铜瓦厢决口改道北…}咸丰五年（1855年），围绕「黄河在铜瓦厢决口改道北流，重塑华北水系与社会经济」的命令、调度、照会或呈报形成独立的中央—地方行动文书链。核心取证：《清文宗实录》咸丰五年相关条；地方志、督抚奏折及地方档案。该项锚定的是命令或决策环节，命令抵达和结果仍须另外核验。来源保留：奏折、上谕、部议、军机档或条约可证明行政动作和机构立场，不自动证明所有地区、群体或对手按其意图行动。
\eventanchor{EQ-1855-YELLOW-RIVER-COURSE-SHIFT-AFTERMATH}{「黄河在铜瓦厢决口改道北流，…}咸丰五年（1855年），「黄河在铜瓦厢决口改道北流，重塑华北水系与社会经济」引发的善后、执行或地方回应构成可由不同材料追踪的后续节点。核心取证：《清文宗实录》咸丰五年相关条；地方志、督抚奏折及地方档案；相关地区的档案、志书、契约、报刊或对方材料。该项锚定的是后续追踪问题，影响范围须按地区与群体分别判断。来源保留：战后、改革后或条约后的记忆、地方记录和官方总结常具有事后选择性；后续影响须按地点、群体和时间分层，不作整体化推断。
\eventanchor{EQ-1856-TIANJING-INCIDENT-PRELUDE}{「天京事变中太平天国领导层发…}咸丰六年（1856年），「天京事变中太平天国领导层发生内讧和大规模杀戮」之前的条件、信息和争议已通过中央与地方材料显现，构成该事件可追踪的前史节点。核心取证：《清文宗实录》咸丰六年相关条；宫中朱批奏折、内阁大库题本与《清史稿》相关志传；同时期地方档案、地方志、日记或对方材料应按具体地区补检。该项锚定的是前史条件与信息背景，不把条件直接写成唯一因果。来源保留：官修编年和地方材料的记录密度与立场并不相同；本行不将条件、传闻、呈报或后见解释直接等同为确定因果。
\eventanchor{EQ-1856-TIANJING-INCIDENT-DECISION}{围绕「天京事变中太平天国领导…}咸丰六年（1856年），围绕「天京事变中太平天国领导层发生内讧和大规模杀戮」的命令、调度、照会或呈报形成独立的中央—地方行动文书链。核心取证：《清文宗实录》咸丰六年相关条；宫中朱批奏折、内阁大库题本与《清史稿》相关志传。该项锚定的是命令或决策环节，命令抵达和结果仍须另外核验。来源保留：奏折、上谕、部议、军机档或条约可证明行政动作和机构立场，不自动证明所有地区、群体或对手按其意图行动。
\eventanchor{EQ-1856-TIANJING-INCIDENT-AFTERMATH}{「天京事变中太平天国领导层发…}咸丰六年（1856年），「天京事变中太平天国领导层发生内讧和大规模杀戮」引发的善后、执行或地方回应构成可由不同材料追踪的后续节点。核心取证：《清文宗实录》咸丰六年相关条；宫中朱批奏折、内阁大库题本与《清史稿》相关志传；相关地区的档案、志书、契约、报刊或对方材料。该项锚定的是后续追踪问题，影响范围须按地区与群体分别判断。来源保留：战后、改革后或条约后的记忆、地方记录和官方总结常具有事后选择性；后续影响须按地点、群体和时间分层，不作整体化推断。
\eventanchor{EQ-1856-SECOND-OPIUM-WAR-BEGIN-PRELUDE}{「亚罗号事件等冲突后英法对清…}咸丰六年（1856年），「亚罗号事件等冲突后英法对清开战，第二次鸦片战争开始」之前的条件、信息和争议已通过中央与地方材料显现，构成该事件可追踪的前史节点。核心取证：《清文宗实录》咸丰六年相关条；军机处档、战事方略及地方奏报；同时期地方档案、地方志、日记或对方材料应按具体地区补检。该项锚定的是前史条件与信息背景，不把条件直接写成唯一因果。来源保留：官修编年和地方材料的记录密度与立场并不相同；本行不将条件、传闻、呈报或后见解释直接等同为确定因果。
\eventanchor{EQ-1856-SECOND-OPIUM-WAR-BEGIN-DECISION}{围绕「亚罗号事件等冲突后英法…}咸丰六年（1856年），围绕「亚罗号事件等冲突后英法对清开战，第二次鸦片战争开始」的命令、调度、照会或呈报形成独立的中央—地方行动文书链。核心取证：《清文宗实录》咸丰六年相关条；军机处档、战事方略及地方奏报。该项锚定的是命令或决策环节，命令抵达和结果仍须另外核验。来源保留：奏折、上谕、部议、军机档或条约可证明行政动作和机构立场，不自动证明所有地区、群体或对手按其意图行动。
\eventanchor{EQ-1856-SECOND-OPIUM-WAR-BEGIN-AFTERMATH}{「亚罗号事件等冲突后英法对清…}咸丰六年（1856年），「亚罗号事件等冲突后英法对清开战，第二次鸦片战争开始」引发的善后、执行或地方回应构成可由不同材料追踪的后续节点。核心取证：《清文宗实录》咸丰六年相关条；军机处档、战事方略及地方奏报；相关地区的档案、志书、契约、报刊或对方材料。该项锚定的是后续追踪问题，影响范围须按地区与群体分别判断。来源保留：战后、改革后或条约后的记忆、地方记录和官方总结常具有事后选择性；后续影响须按地点、群体和时间分层，不作整体化推断。
\eventanchor{EQ-1857-GUANGZHOU-FALL-PRELUDE}{「英法联军攻占广州，叶名琛被…}咸丰七年（1857年），「英法联军攻占广州，叶名琛被俘」之前的条件、信息和争议已通过中央与地方材料显现，构成该事件可追踪的前史节点。核心取证：《清文宗实录》咸丰七年相关条；军机处档、战事方略及地方奏报；同时期地方档案、地方志、日记或对方材料应按具体地区补检。该项锚定的是前史条件与信息背景，不把条件直接写成唯一因果。来源保留：官修编年和地方材料的记录密度与立场并不相同；本行不将条件、传闻、呈报或后见解释直接等同为确定因果。
\eventanchor{EQ-1857-GUANGZHOU-FALL-DECISION}{围绕「英法联军攻占广州，叶名…}咸丰七年（1857年），围绕「英法联军攻占广州，叶名琛被俘」的命令、调度、照会或呈报形成独立的中央—地方行动文书链。核心取证：《清文宗实录》咸丰七年相关条；军机处档、战事方略及地方奏报。该项锚定的是命令或决策环节，命令抵达和结果仍须另外核验。来源保留：奏折、上谕、部议、军机档或条约可证明行政动作和机构立场，不自动证明所有地区、群体或对手按其意图行动。
\eventanchor{EQ-1857-GUANGZHOU-FALL-AFTERMATH}{「英法联军攻占广州，叶名琛被…}咸丰七年（1857年），「英法联军攻占广州，叶名琛被俘」引发的善后、执行或地方回应构成可由不同材料追踪的后续节点。核心取证：《清文宗实录》咸丰七年相关条；军机处档、战事方略及地方奏报；相关地区的档案、志书、契约、报刊或对方材料。该项锚定的是后续追踪问题，影响范围须按地区与群体分别判断。来源保留：战后、改革后或条约后的记忆、地方记录和官方总结常具有事后选择性；后续影响须按地点、群体和时间分层，不作整体化推断。
\eventanchor{EQ-1859-TAKU-BATTLE-PRELUDE}{「大沽口战事使英法换约受阻，…}咸丰九年（1859年），「大沽口战事使英法换约受阻，战争再度升级」之前的条件、信息和争议已通过中央与地方材料显现，构成该事件可追踪的前史节点。核心取证：《清文宗实录》咸丰九年相关条；军机处档、战事方略及地方奏报；同时期地方档案、地方志、日记或对方材料应按具体地区补检。该项锚定的是前史条件与信息背景，不把条件直接写成唯一因果。来源保留：官修编年和地方材料的记录密度与立场并不相同；本行不将条件、传闻、呈报或后见解释直接等同为确定因果。
\eventanchor{EQ-1859-TAKU-BATTLE-DECISION}{围绕「大沽口战事使英法换约受…}咸丰九年（1859年），围绕「大沽口战事使英法换约受阻，战争再度升级」的命令、调度、照会或呈报形成独立的中央—地方行动文书链。核心取证：《清文宗实录》咸丰九年相关条；军机处档、战事方略及地方奏报。该项锚定的是命令或决策环节，命令抵达和结果仍须另外核验。来源保留：奏折、上谕、部议、军机档或条约可证明行政动作和机构立场，不自动证明所有地区、群体或对手按其意图行动。
\eventanchor{EQ-1859-TAKU-BATTLE-AFTERMATH}{「大沽口战事使英法换约受阻，…}咸丰九年（1859年），「大沽口战事使英法换约受阻，战争再度升级」引发的善后、执行或地方回应构成可由不同材料追踪的后续节点。核心取证：《清文宗实录》咸丰九年相关条；军机处档、战事方略及地方奏报；相关地区的档案、志书、契约、报刊或对方材料。该项锚定的是后续追踪问题，影响范围须按地区与群体分别判断。来源保留：战后、改革后或条约后的记忆、地方记录和官方总结常具有事后选择性；后续影响须按地点、群体和时间分层，不作整体化推断。
\subsection{1860—1869}
\eventanchor{EQ-1860-BEIJING-CAMPAIGN-PRELUDE}{「英法联军进逼北京，圆明园遭…}咸丰十年（1860年），「英法联军进逼北京，圆明园遭焚掠，咸丰帝出走热河」之前的条件、信息和争议已通过中央与地方材料显现，构成该事件可追踪的前史节点。核心取证：《清文宗实录》咸丰十年相关条；军机处档、战事方略及地方奏报；同时期地方档案、地方志、日记或对方材料应按具体地区补检。该项锚定的是前史条件与信息背景，不把条件直接写成唯一因果。来源保留：官修编年和地方材料的记录密度与立场并不相同；本行不将条件、传闻、呈报或后见解释直接等同为确定因果。
\eventanchor{EQ-1860-BEIJING-CAMPAIGN-DECISION}{围绕「英法联军进逼北京，圆明…}咸丰十年（1860年），围绕「英法联军进逼北京，圆明园遭焚掠，咸丰帝出走热河」的命令、调度、照会或呈报形成独立的中央—地方行动文书链。核心取证：《清文宗实录》咸丰十年相关条；军机处档、战事方略及地方奏报。该项锚定的是命令或决策环节，命令抵达和结果仍须另外核验。来源保留：奏折、上谕、部议、军机档或条约可证明行政动作和机构立场，不自动证明所有地区、群体或对手按其意图行动。
\eventanchor{EQ-1860-BEIJING-CAMPAIGN-AFTERMATH}{「英法联军进逼北京，圆明园遭…}咸丰十年（1860年），「英法联军进逼北京，圆明园遭焚掠，咸丰帝出走热河」引发的善后、执行或地方回应构成可由不同材料追踪的后续节点。核心取证：《清文宗实录》咸丰十年相关条；军机处档、战事方略及地方奏报；相关地区的档案、志书、契约、报刊或对方材料。该项锚定的是后续追踪问题，影响范围须按地区与群体分别判断。来源保留：战后、改革后或条约后的记忆、地方记录和官方总结常具有事后选择性；后续影响须按地点、群体和时间分层，不作整体化推断。
\eventanchor{EQ-1860-CONVENTION-OF-PEKING-PRELUDE}{「《北京条约》及相关中俄条约…}咸丰十年（1860年），「《北京条约》及相关中俄条约确认和扩大列强权益并涉及东北边界」之前的条件、信息和争议已通过中央与地方材料显现，构成该事件可追踪的前史节点。核心取证：《清文宗实录》咸丰十年相关条；《筹办夷务始末》相关年卷与中外照会、条约文本；同时期地方档案、地方志、日记或对方材料应按具体地区补检。该项锚定的是前史条件与信息背景，不把条件直接写成唯一因果。来源保留：官修编年和地方材料的记录密度与立场并不相同；本行不将条件、传闻、呈报或后见解释直接等同为确定因果。
\eventanchor{EQ-1860-CONVENTION-OF-PEKING-DECISION}{围绕「《北京条约》及相关中俄…}咸丰十年（1860年），围绕「《北京条约》及相关中俄条约确认和扩大列强权益并涉及东北边界」的命令、调度、照会或呈报形成独立的中央—地方行动文书链。核心取证：《清文宗实录》咸丰十年相关条；《筹办夷务始末》相关年卷与中外照会、条约文本。该项锚定的是命令或决策环节，命令抵达和结果仍须另外核验。来源保留：奏折、上谕、部议、军机档或条约可证明行政动作和机构立场，不自动证明所有地区、群体或对手按其意图行动。
\eventanchor{EQ-1860-CONVENTION-OF-PEKING-AFTERMATH}{「《北京条约》及相关中俄条约…}咸丰十年（1860年），「《北京条约》及相关中俄条约确认和扩大列强权益并涉及东北边界」引发的善后、执行或地方回应构成可由不同材料追踪的后续节点。核心取证：《清文宗实录》咸丰十年相关条；《筹办夷务始末》相关年卷与中外照会、条约文本；相关地区的档案、志书、契约、报刊或对方材料。该项锚定的是后续追踪问题，影响范围须按地区与群体分别判断。来源保留：战后、改革后或条约后的记忆、地方记录和官方总结常具有事后选择性；后续影响须按地点、群体和时间分层，不作整体化推断。
\eventanchor{EQ-1861-XINYOU-COUP-PRELUDE}{「慈禧、慈安与恭亲王发动辛酉…}咸丰十一年（1861年），「慈禧、慈安与恭亲王发动辛酉政变，肃顺等顾命大臣失势」之前的条件、信息和争议已通过中央与地方材料显现，构成该事件可追踪的前史节点。核心取证：《清文宗实录》咸丰十一年相关条；宫中朱批奏折、内阁大库题本与《清史稿》相关志传；同时期地方档案、地方志、日记或对方材料应按具体地区补检。该项锚定的是前史条件与信息背景，不把条件直接写成唯一因果。来源保留：官修编年和地方材料的记录密度与立场并不相同；本行不将条件、传闻、呈报或后见解释直接等同为确定因果。
\eventanchor{EQ-1861-XINYOU-COUP-DECISION}{围绕「慈禧、慈安与恭亲王发动…}咸丰十一年（1861年），围绕「慈禧、慈安与恭亲王发动辛酉政变，肃顺等顾命大臣失势」的命令、调度、照会或呈报形成独立的中央—地方行动文书链。核心取证：《清文宗实录》咸丰十一年相关条；宫中朱批奏折、内阁大库题本与《清史稿》相关志传。该项锚定的是命令或决策环节，命令抵达和结果仍须另外核验。来源保留：奏折、上谕、部议、军机档或条约可证明行政动作和机构立场，不自动证明所有地区、群体或对手按其意图行动。
\eventanchor{EQ-1861-XINYOU-COUP-AFTERMATH}{「慈禧、慈安与恭亲王发动辛酉…}咸丰十一年（1861年），「慈禧、慈安与恭亲王发动辛酉政变，肃顺等顾命大臣失势」引发的善后、执行或地方回应构成可由不同材料追踪的后续节点。核心取证：《清文宗实录》咸丰十一年相关条；宫中朱批奏折、内阁大库题本与《清史稿》相关志传；相关地区的档案、志书、契约、报刊或对方材料。该项锚定的是后续追踪问题，影响范围须按地区与群体分别判断。来源保留：战后、改革后或条约后的记忆、地方记录和官方总结常具有事后选择性；后续影响须按地点、群体和时间分层，不作整体化推断。
\eventanchor{EQ-1861-ZONGLI-YAMEN-PRELUDE}{「总理各国事务衙门设立，清廷…}咸丰十一年（1861年），「总理各国事务衙门设立，清廷开始形成常设近代外交处理机构」之前的条件、信息和争议已通过中央与地方材料显现，构成该事件可追踪的前史节点。核心取证：《清文宗实录》咸丰十一年相关条；《筹办夷务始末》相关年卷与中外照会、条约文本；同时期地方档案、地方志、日记或对方材料应按具体地区补检。该项锚定的是前史条件与信息背景，不把条件直接写成唯一因果。来源保留：官修编年和地方材料的记录密度与立场并不相同；本行不将条件、传闻、呈报或后见解释直接等同为确定因果。
\eventanchor{EQ-1861-ZONGLI-YAMEN-DECISION}{围绕「总理各国事务衙门设立，…}咸丰十一年（1861年），围绕「总理各国事务衙门设立，清廷开始形成常设近代外交处理机构」的命令、调度、照会或呈报形成独立的中央—地方行动文书链。核心取证：《清文宗实录》咸丰十一年相关条；《筹办夷务始末》相关年卷与中外照会、条约文本。该项锚定的是命令或决策环节，命令抵达和结果仍须另外核验。来源保留：奏折、上谕、部议、军机档或条约可证明行政动作和机构立场，不自动证明所有地区、群体或对手按其意图行动。
\eventanchor{EQ-1861-ZONGLI-YAMEN-AFTERMATH}{「总理各国事务衙门设立，清廷…}咸丰十一年（1861年），「总理各国事务衙门设立，清廷开始形成常设近代外交处理机构」引发的善后、执行或地方回应构成可由不同材料追踪的后续节点。核心取证：《清文宗实录》咸丰十一年相关条；《筹办夷务始末》相关年卷与中外照会、条约文本；相关地区的档案、志书、契约、报刊或对方材料。该项锚定的是后续追踪问题，影响范围须按地区与群体分别判断。来源保留：战后、改革后或条约后的记忆、地方记录和官方总结常具有事后选择性；后续影响须按地点、群体和时间分层，不作整体化推断。
\eventanchor{EQ-1862-SELF-STRENGTHENING-ORIGINS-PRELUDE}{「洋务实践在军械、翻译、外交…}同治元年（1862年），「洋务实践在军械、翻译、外交和地方军政中展开」之前的条件、信息和争议已通过中央与地方材料显现，构成该事件可追踪的前史节点。核心取证：《清穆宗实录》同治元年相关条；宫中朱批奏折、内阁大库题本与《清史稿》相关志传；同时期地方档案、地方志、日记或对方材料应按具体地区补检。该项锚定的是前史条件与信息背景，不把条件直接写成唯一因果。来源保留：官修编年和地方材料的记录密度与立场并不相同；本行不将条件、传闻、呈报或后见解释直接等同为确定因果。
\eventanchor{EQ-1862-SELF-STRENGTHENING-ORIGINS-DECISION}{围绕「洋务实践在军械、翻译、…}同治元年（1862年），围绕「洋务实践在军械、翻译、外交和地方军政中展开」的命令、调度、照会或呈报形成独立的中央—地方行动文书链。核心取证：《清穆宗实录》同治元年相关条；宫中朱批奏折、内阁大库题本与《清史稿》相关志传。该项锚定的是命令或决策环节，命令抵达和结果仍须另外核验。来源保留：奏折、上谕、部议、军机档或条约可证明行政动作和机构立场，不自动证明所有地区、群体或对手按其意图行动。
\eventanchor{EQ-1862-SELF-STRENGTHENING-ORIGINS-AFTERMATH}{「洋务实践在军械、翻译、外交…}同治元年（1862年），「洋务实践在军械、翻译、外交和地方军政中展开」引发的善后、执行或地方回应构成可由不同材料追踪的后续节点。核心取证：《清穆宗实录》同治元年相关条；宫中朱批奏折、内阁大库题本与《清史稿》相关志传；相关地区的档案、志书、契约、报刊或对方材料。该项锚定的是后续追踪问题，影响范围须按地区与群体分别判断。来源保留：战后、改革后或条约后的记忆、地方记录和官方总结常具有事后选择性；后续影响须按地点、群体和时间分层，不作整体化推断。
\eventanchor{EQ-1863-EVER-VICTORIOUS-ARMY-PRELUDE}{「常胜军在江南战场参与清军作…}同治二年（1863年），「常胜军在江南战场参与清军作战，戈登接掌后扩大影响」之前的条件、信息和争议已通过中央与地方材料显现，构成该事件可追踪的前史节点。核心取证：《清穆宗实录》同治二年相关条；军机处档、战事方略及地方奏报；同时期地方档案、地方志、日记或对方材料应按具体地区补检。该项锚定的是前史条件与信息背景，不把条件直接写成唯一因果。来源保留：官修编年和地方材料的记录密度与立场并不相同；本行不将条件、传闻、呈报或后见解释直接等同为确定因果。
\eventanchor{EQ-1863-EVER-VICTORIOUS-ARMY-DECISION}{围绕「常胜军在江南战场参与清…}同治二年（1863年），围绕「常胜军在江南战场参与清军作战，戈登接掌后扩大影响」的命令、调度、照会或呈报形成独立的中央—地方行动文书链。核心取证：《清穆宗实录》同治二年相关条；军机处档、战事方略及地方奏报。该项锚定的是命令或决策环节，命令抵达和结果仍须另外核验。来源保留：奏折、上谕、部议、军机档或条约可证明行政动作和机构立场，不自动证明所有地区、群体或对手按其意图行动。
\eventanchor{EQ-1863-EVER-VICTORIOUS-ARMY-AFTERMATH}{「常胜军在江南战场参与清军作…}同治二年（1863年），「常胜军在江南战场参与清军作战，戈登接掌后扩大影响」引发的善后、执行或地方回应构成可由不同材料追踪的后续节点。核心取证：《清穆宗实录》同治二年相关条；军机处档、战事方略及地方奏报；相关地区的档案、志书、契约、报刊或对方材料。该项锚定的是后续追踪问题，影响范围须按地区与群体分别判断。来源保留：战后、改革后或条约后的记忆、地方记录和官方总结常具有事后选择性；后续影响须按地点、群体和时间分层，不作整体化推断。
\eventanchor{EQ-1863-NIAN-WAR-MOBILIZATION-PRELUDE}{「清廷调集湘淮军和地方力量应…}同治二年（1863年），「清廷调集湘淮军和地方力量应对捻军，华北战事加剧」之前的条件、信息和争议已通过中央与地方材料显现，构成该事件可追踪的前史节点。核心取证：《清穆宗实录》同治二年相关条；军机处档、战事方略及地方奏报；同时期地方档案、地方志、日记或对方材料应按具体地区补检。该项锚定的是前史条件与信息背景，不把条件直接写成唯一因果。来源保留：官修编年和地方材料的记录密度与立场并不相同；本行不将条件、传闻、呈报或后见解释直接等同为确定因果。
\eventanchor{EQ-1863-NIAN-WAR-MOBILIZATION-DECISION}{围绕「清廷调集湘淮军和地方力…}同治二年（1863年），围绕「清廷调集湘淮军和地方力量应对捻军，华北战事加剧」的命令、调度、照会或呈报形成独立的中央—地方行动文书链。核心取证：《清穆宗实录》同治二年相关条；军机处档、战事方略及地方奏报。该项锚定的是命令或决策环节，命令抵达和结果仍须另外核验。来源保留：奏折、上谕、部议、军机档或条约可证明行政动作和机构立场，不自动证明所有地区、群体或对手按其意图行动。
\eventanchor{EQ-1863-NIAN-WAR-MOBILIZATION-AFTERMATH}{「清廷调集湘淮军和地方力量应…}同治二年（1863年），「清廷调集湘淮军和地方力量应对捻军，华北战事加剧」引发的善后、执行或地方回应构成可由不同材料追踪的后续节点。核心取证：《清穆宗实录》同治二年相关条；军机处档、战事方略及地方奏报；相关地区的档案、志书、契约、报刊或对方材料。该项锚定的是后续追踪问题，影响范围须按地区与群体分别判断。来源保留：战后、改革后或条约后的记忆、地方记录和官方总结常具有事后选择性；后续影响须按地点、群体和时间分层，不作整体化推断。
\eventanchor{EQ-1864-TIANJING-FALL-PRELUDE}{「湘军攻陷天京，太平天国核心…}同治三年（1864年），「湘军攻陷天京，太平天国核心政权覆亡」之前的条件、信息和争议已通过中央与地方材料显现，构成该事件可追踪的前史节点。核心取证：《清穆宗实录》同治三年相关条；军机处档、战事方略及地方奏报；同时期地方档案、地方志、日记或对方材料应按具体地区补检。该项锚定的是前史条件与信息背景，不把条件直接写成唯一因果。来源保留：官修编年和地方材料的记录密度与立场并不相同；本行不将条件、传闻、呈报或后见解释直接等同为确定因果。
\eventanchor{EQ-1864-TIANJING-FALL-DECISION}{围绕「湘军攻陷天京，太平天国…}同治三年（1864年），围绕「湘军攻陷天京，太平天国核心政权覆亡」的命令、调度、照会或呈报形成独立的中央—地方行动文书链。核心取证：《清穆宗实录》同治三年相关条；军机处档、战事方略及地方奏报。该项锚定的是命令或决策环节，命令抵达和结果仍须另外核验。来源保留：奏折、上谕、部议、军机档或条约可证明行政动作和机构立场，不自动证明所有地区、群体或对手按其意图行动。
\eventanchor{EQ-1864-TIANJING-FALL-AFTERMATH}{「湘军攻陷天京，太平天国核心…}同治三年（1864年），「湘军攻陷天京，太平天国核心政权覆亡」引发的善后、执行或地方回应构成可由不同材料追踪的后续节点。核心取证：《清穆宗实录》同治三年相关条；军机处档、战事方略及地方奏报；相关地区的档案、志书、契约、报刊或对方材料。该项锚定的是后续追踪问题，影响范围须按地区与群体分别判断。来源保留：战后、改革后或条约后的记忆、地方记录和官方总结常具有事后选择性；后续影响须按地点、群体和时间分层，不作整体化推断。
\eventanchor{EQ-1865-DUNGAN-REVOLT-PRELUDE}{「西北回民起事在陕西、甘肃等…}同治四年（1865年），「西北回民起事在陕西、甘肃等地扩展，形成长期战争」之前的条件、信息和争议已通过中央与地方材料显现，构成该事件可追踪的前史节点。核心取证：《清穆宗实录》同治四年相关条；军机处档、战事方略及地方奏报；同时期地方档案、地方志、日记或对方材料应按具体地区补检。该项锚定的是前史条件与信息背景，不把条件直接写成唯一因果。来源保留：官修编年和地方材料的记录密度与立场并不相同；本行不将条件、传闻、呈报或后见解释直接等同为确定因果。
\eventanchor{EQ-1865-DUNGAN-REVOLT-DECISION}{围绕「西北回民起事在陕西、甘…}同治四年（1865年），围绕「西北回民起事在陕西、甘肃等地扩展，形成长期战争」的命令、调度、照会或呈报形成独立的中央—地方行动文书链。核心取证：《清穆宗实录》同治四年相关条；军机处档、战事方略及地方奏报。该项锚定的是命令或决策环节，命令抵达和结果仍须另外核验。来源保留：奏折、上谕、部议、军机档或条约可证明行政动作和机构立场，不自动证明所有地区、群体或对手按其意图行动。
\eventanchor{EQ-1865-DUNGAN-REVOLT-AFTERMATH}{「西北回民起事在陕西、甘肃等…}同治四年（1865年），「西北回民起事在陕西、甘肃等地扩展，形成长期战争」引发的善后、执行或地方回应构成可由不同材料追踪的后续节点。核心取证：《清穆宗实录》同治四年相关条；军机处档、战事方略及地方奏报；相关地区的档案、志书、契约、报刊或对方材料。该项锚定的是后续追踪问题，影响范围须按地区与群体分别判断。来源保留：战后、改革后或条约后的记忆、地方记录和官方总结常具有事后选择性；后续影响须按地点、群体和时间分层，不作整体化推断。
\eventanchor{EQ-1865-NIAN-LEADER-DEATH-PRELUDE}{「张乐行等捻军领导人先后失势…}同治四年（1865年），「张乐行等捻军领导人先后失势或被杀，捻军仍持续活动」之前的条件、信息和争议已通过中央与地方材料显现，构成该事件可追踪的前史节点。核心取证：《清穆宗实录》同治四年相关条；军机处档、战事方略及地方奏报；同时期地方档案、地方志、日记或对方材料应按具体地区补检。该项锚定的是前史条件与信息背景，不把条件直接写成唯一因果。来源保留：官修编年和地方材料的记录密度与立场并不相同；本行不将条件、传闻、呈报或后见解释直接等同为确定因果。
\eventanchor{EQ-1865-NIAN-LEADER-DEATH-DECISION}{围绕「张乐行等捻军领导人先后…}同治四年（1865年），围绕「张乐行等捻军领导人先后失势或被杀，捻军仍持续活动」的命令、调度、照会或呈报形成独立的中央—地方行动文书链。核心取证：《清穆宗实录》同治四年相关条；军机处档、战事方略及地方奏报。该项锚定的是命令或决策环节，命令抵达和结果仍须另外核验。来源保留：奏折、上谕、部议、军机档或条约可证明行政动作和机构立场，不自动证明所有地区、群体或对手按其意图行动。
\eventanchor{EQ-1865-NIAN-LEADER-DEATH-AFTERMATH}{「张乐行等捻军领导人先后失势…}同治四年（1865年），「张乐行等捻军领导人先后失势或被杀，捻军仍持续活动」引发的善后、执行或地方回应构成可由不同材料追踪的后续节点。核心取证：《清穆宗实录》同治四年相关条；军机处档、战事方略及地方奏报；相关地区的档案、志书、契约、报刊或对方材料。该项锚定的是后续追踪问题，影响范围须按地区与群体分别判断。来源保留：战后、改革后或条约后的记忆、地方记录和官方总结常具有事后选择性；后续影响须按地点、群体和时间分层，不作整体化推断。
\eventanchor{EQ-1866-FUJIAN-SHIPYARD-PRELUDE}{「福州船政局创办，成为洋务造…}同治五年（1866年），「福州船政局创办，成为洋务造船、教育与技术翻译的重要项目」之前的条件、信息和争议已通过中央与地方材料显现，构成该事件可追踪的前史节点。核心取证：《清穆宗实录》同治五年相关条；户部奏销、盐政、河工或海关档案；同时期地方档案、地方志、日记或对方材料应按具体地区补检。该项锚定的是前史条件与信息背景，不把条件直接写成唯一因果。来源保留：官修编年和地方材料的记录密度与立场并不相同；本行不将条件、传闻、呈报或后见解释直接等同为确定因果。
\eventanchor{EQ-1866-FUJIAN-SHIPYARD-DECISION}{围绕「福州船政局创办，成为洋…}同治五年（1866年），围绕「福州船政局创办，成为洋务造船、教育与技术翻译的重要项目」的命令、调度、照会或呈报形成独立的中央—地方行动文书链。核心取证：《清穆宗实录》同治五年相关条；户部奏销、盐政、河工或海关档案。该项锚定的是命令或决策环节，命令抵达和结果仍须另外核验。来源保留：奏折、上谕、部议、军机档或条约可证明行政动作和机构立场，不自动证明所有地区、群体或对手按其意图行动。
\eventanchor{EQ-1866-FUJIAN-SHIPYARD-AFTERMATH}{「福州船政局创办，成为洋务造…}同治五年（1866年），「福州船政局创办，成为洋务造船、教育与技术翻译的重要项目」引发的善后、执行或地方回应构成可由不同材料追踪的后续节点。核心取证：《清穆宗实录》同治五年相关条；户部奏销、盐政、河工或海关档案；相关地区的档案、志书、契约、报刊或对方材料。该项锚定的是后续追踪问题，影响范围须按地区与群体分别判断。来源保留：战后、改革后或条约后的记忆、地方记录和官方总结常具有事后选择性；后续影响须按地点、群体和时间分层，不作整体化推断。
\eventanchor{EQ-1866-NINGBO-SHANGHAI-INDUSTRY-PRELUDE}{「江南制造局等军工项目在战后…}同治五年（1866年），「江南制造局等军工项目在战后江南环境中继续发展」之前的条件、信息和争议已通过中央与地方材料显现，构成该事件可追踪的前史节点。核心取证：《清穆宗实录》同治五年相关条；户部奏销、盐政、河工或海关档案；同时期地方档案、地方志、日记或对方材料应按具体地区补检。该项锚定的是前史条件与信息背景，不把条件直接写成唯一因果。来源保留：官修编年和地方材料的记录密度与立场并不相同；本行不将条件、传闻、呈报或后见解释直接等同为确定因果。
\eventanchor{EQ-1866-NINGBO-SHANGHAI-INDUSTRY-DECISION}{围绕「江南制造局等军工项目在…}同治五年（1866年），围绕「江南制造局等军工项目在战后江南环境中继续发展」的命令、调度、照会或呈报形成独立的中央—地方行动文书链。核心取证：《清穆宗实录》同治五年相关条；户部奏销、盐政、河工或海关档案。该项锚定的是命令或决策环节，命令抵达和结果仍须另外核验。来源保留：奏折、上谕、部议、军机档或条约可证明行政动作和机构立场，不自动证明所有地区、群体或对手按其意图行动。
\eventanchor{EQ-1866-NINGBO-SHANGHAI-INDUSTRY-AFTERMATH}{「江南制造局等军工项目在战后…}同治五年（1866年），「江南制造局等军工项目在战后江南环境中继续发展」引发的善后、执行或地方回应构成可由不同材料追踪的后续节点。核心取证：《清穆宗实录》同治五年相关条；户部奏销、盐政、河工或海关档案；相关地区的档案、志书、契约、报刊或对方材料。该项锚定的是后续追踪问题，影响范围须按地区与群体分别判断。来源保留：战后、改革后或条约后的记忆、地方记录和官方总结常具有事后选择性；后续影响须按地点、群体和时间分层，不作整体化推断。
\eventanchor{EQ-1868-NIAN-SUPPRESSION-PRELUDE}{「东捻军等主力被击败，捻军大…}同治七年（1868年），「东捻军等主力被击败，捻军大规模战争告终」之前的条件、信息和争议已通过中央与地方材料显现，构成该事件可追踪的前史节点。核心取证：《清穆宗实录》同治七年相关条；军机处档、战事方略及地方奏报；同时期地方档案、地方志、日记或对方材料应按具体地区补检。该项锚定的是前史条件与信息背景，不把条件直接写成唯一因果。来源保留：官修编年和地方材料的记录密度与立场并不相同；本行不将条件、传闻、呈报或后见解释直接等同为确定因果。
\eventanchor{EQ-1868-NIAN-SUPPRESSION-DECISION}{围绕「东捻军等主力被击败，捻…}同治七年（1868年），围绕「东捻军等主力被击败，捻军大规模战争告终」的命令、调度、照会或呈报形成独立的中央—地方行动文书链。核心取证：《清穆宗实录》同治七年相关条；军机处档、战事方略及地方奏报。该项锚定的是命令或决策环节，命令抵达和结果仍须另外核验。来源保留：奏折、上谕、部议、军机档或条约可证明行政动作和机构立场，不自动证明所有地区、群体或对手按其意图行动。
\eventanchor{EQ-1868-NIAN-SUPPRESSION-AFTERMATH}{「东捻军等主力被击败，捻军大…}同治七年（1868年），「东捻军等主力被击败，捻军大规模战争告终」引发的善后、执行或地方回应构成可由不同材料追踪的后续节点。核心取证：《清穆宗实录》同治七年相关条；军机处档、战事方略及地方奏报；相关地区的档案、志书、契约、报刊或对方材料。该项锚定的是后续追踪问题，影响范围须按地区与群体分别判断。来源保留：战后、改革后或条约后的记忆、地方记录和官方总结常具有事后选择性；后续影响须按地点、群体和时间分层，不作整体化推断。
\eventanchor{EQ-1869-YANGZHOU-SALT-CRISIS-PRELUDE}{「两淮盐政改革、商欠与财政重…}同治八年（1869年），「两淮盐政改革、商欠与财政重整成为中兴时期的重要议题」之前的条件、信息和争议已通过中央与地方材料显现，构成该事件可追踪的前史节点。核心取证：《清穆宗实录》同治八年相关条；户部奏销、盐政、河工或海关档案；同时期地方档案、地方志、日记或对方材料应按具体地区补检。该项锚定的是前史条件与信息背景，不把条件直接写成唯一因果。来源保留：官修编年和地方材料的记录密度与立场并不相同；本行不将条件、传闻、呈报或后见解释直接等同为确定因果。
\eventanchor{EQ-1869-YANGZHOU-SALT-CRISIS-DECISION}{围绕「两淮盐政改革、商欠与财…}同治八年（1869年），围绕「两淮盐政改革、商欠与财政重整成为中兴时期的重要议题」的命令、调度、照会或呈报形成独立的中央—地方行动文书链。核心取证：《清穆宗实录》同治八年相关条；户部奏销、盐政、河工或海关档案。该项锚定的是命令或决策环节，命令抵达和结果仍须另外核验。来源保留：奏折、上谕、部议、军机档或条约可证明行政动作和机构立场，不自动证明所有地区、群体或对手按其意图行动。
\eventanchor{EQ-1869-YANGZHOU-SALT-CRISIS-AFTERMATH}{「两淮盐政改革、商欠与财政重…}同治八年（1869年），「两淮盐政改革、商欠与财政重整成为中兴时期的重要议题」引发的善后、执行或地方回应构成可由不同材料追踪的后续节点。核心取证：《清穆宗实录》同治八年相关条；户部奏销、盐政、河工或海关档案；相关地区的档案、志书、契约、报刊或对方材料。该项锚定的是后续追踪问题，影响范围须按地区与群体分别判断。来源保留：战后、改革后或条约后的记忆、地方记录和官方总结常具有事后选择性；后续影响须按地点、群体和时间分层，不作整体化推断。
\subsection{1870—1879}
\eventanchor{EQ-1870-TIANJIN-MASSACRE-PRELUDE}{「天津教案发生，法国领事、修…}同治九年（1870年），「天津教案发生，法国领事、修女和中国民众死伤，引发外交危机」之前的条件、信息和争议已通过中央与地方材料显现，构成该事件可追踪的前史节点。核心取证：《清穆宗实录》同治九年相关条；《筹办夷务始末》相关年卷与中外照会、条约文本；同时期地方档案、地方志、日记或对方材料应按具体地区补检。该项锚定的是前史条件与信息背景，不把条件直接写成唯一因果。来源保留：官修编年和地方材料的记录密度与立场并不相同；本行不将条件、传闻、呈报或后见解释直接等同为确定因果。
\eventanchor{EQ-1870-TIANJIN-MASSACRE-DECISION}{围绕「天津教案发生，法国领事…}同治九年（1870年），围绕「天津教案发生，法国领事、修女和中国民众死伤，引发外交危机」的命令、调度、照会或呈报形成独立的中央—地方行动文书链。核心取证：《清穆宗实录》同治九年相关条；《筹办夷务始末》相关年卷与中外照会、条约文本。该项锚定的是命令或决策环节，命令抵达和结果仍须另外核验。来源保留：奏折、上谕、部议、军机档或条约可证明行政动作和机构立场，不自动证明所有地区、群体或对手按其意图行动。
\eventanchor{EQ-1870-TIANJIN-MASSACRE-AFTERMATH}{「天津教案发生，法国领事、修…}同治九年（1870年），「天津教案发生，法国领事、修女和中国民众死伤，引发外交危机」引发的善后、执行或地方回应构成可由不同材料追踪的后续节点。核心取证：《清穆宗实录》同治九年相关条；《筹办夷务始末》相关年卷与中外照会、条约文本；相关地区的档案、志书、契约、报刊或对方材料。该项锚定的是后续追踪问题，影响范围须按地区与群体分别判断。来源保留：战后、改革后或条约后的记忆、地方记录和官方总结常具有事后选择性；后续影响须按地点、群体和时间分层，不作整体化推断。
\eventanchor{EQ-1871-MUDAN-INCIDENT-PRELUDE}{「琉球船民在台湾牡丹社一带遇…}同治十年（1871年），「琉球船民在台湾牡丹社一带遇害，事件成为日本对台政策借口」之前的条件、信息和争议已通过中央与地方材料显现，构成该事件可追踪的前史节点。核心取证：《清穆宗实录》同治十年相关条；《筹办夷务始末》相关年卷与中外照会、条约文本；同时期地方档案、地方志、日记或对方材料应按具体地区补检。该项锚定的是前史条件与信息背景，不把条件直接写成唯一因果。来源保留：官修编年和地方材料的记录密度与立场并不相同；本行不将条件、传闻、呈报或后见解释直接等同为确定因果。
\eventanchor{EQ-1871-MUDAN-INCIDENT-DECISION}{围绕「琉球船民在台湾牡丹社一…}同治十年（1871年），围绕「琉球船民在台湾牡丹社一带遇害，事件成为日本对台政策借口」的命令、调度、照会或呈报形成独立的中央—地方行动文书链。核心取证：《清穆宗实录》同治十年相关条；《筹办夷务始末》相关年卷与中外照会、条约文本。该项锚定的是命令或决策环节，命令抵达和结果仍须另外核验。来源保留：奏折、上谕、部议、军机档或条约可证明行政动作和机构立场，不自动证明所有地区、群体或对手按其意图行动。
\eventanchor{EQ-1871-MUDAN-INCIDENT-AFTERMATH}{「琉球船民在台湾牡丹社一带遇…}同治十年（1871年），「琉球船民在台湾牡丹社一带遇害，事件成为日本对台政策借口」引发的善后、执行或地方回应构成可由不同材料追踪的后续节点。核心取证：《清穆宗实录》同治十年相关条；《筹办夷务始末》相关年卷与中外照会、条约文本；相关地区的档案、志书、契约、报刊或对方材料。该项锚定的是后续追踪问题，影响范围须按地区与群体分别判断。来源保留：战后、改革后或条约后的记忆、地方记录和官方总结常具有事后选择性；后续影响须按地点、群体和时间分层，不作整体化推断。
\eventanchor{EQ-1871-QING-STUDENTS-ABROAD-PRELUDE}{「清廷派遣幼童赴美留学的计划…}同治十年（1871年），「清廷派遣幼童赴美留学的计划启动，洋务教育进入海外培训阶段」之前的条件、信息和争议已通过中央与地方材料显现，构成该事件可追踪的前史节点。核心取证：《清穆宗实录》同治十年相关条；内阁档、文集或报刊相关记录；同时期地方档案、地方志、日记或对方材料应按具体地区补检。该项锚定的是前史条件与信息背景，不把条件直接写成唯一因果。来源保留：官修编年和地方材料的记录密度与立场并不相同；本行不将条件、传闻、呈报或后见解释直接等同为确定因果。
\eventanchor{EQ-1871-QING-STUDENTS-ABROAD-DECISION}{围绕「清廷派遣幼童赴美留学的…}同治十年（1871年），围绕「清廷派遣幼童赴美留学的计划启动，洋务教育进入海外培训阶段」的命令、调度、照会或呈报形成独立的中央—地方行动文书链。核心取证：《清穆宗实录》同治十年相关条；内阁档、文集或报刊相关记录。该项锚定的是命令或决策环节，命令抵达和结果仍须另外核验。来源保留：奏折、上谕、部议、军机档或条约可证明行政动作和机构立场，不自动证明所有地区、群体或对手按其意图行动。
\eventanchor{EQ-1871-QING-STUDENTS-ABROAD-AFTERMATH}{「清廷派遣幼童赴美留学的计划…}同治十年（1871年），「清廷派遣幼童赴美留学的计划启动，洋务教育进入海外培训阶段」引发的善后、执行或地方回应构成可由不同材料追踪的后续节点。核心取证：《清穆宗实录》同治十年相关条；内阁档、文集或报刊相关记录；相关地区的档案、志书、契约、报刊或对方材料。该项锚定的是后续追踪问题，影响范围须按地区与群体分别判断。来源保留：战后、改革后或条约后的记忆、地方记录和官方总结常具有事后选择性；后续影响须按地点、群体和时间分层，不作整体化推断。
\eventanchor{EQ-1872-CHINA-MERCHANTS-STEAM-NAVIGATION-PRELUDE}{「轮船招商局创办，试图以官督…}同治十一年（1872年），「轮船招商局创办，试图以官督商办方式发展近代航运」之前的条件、信息和争议已通过中央与地方材料显现，构成该事件可追踪的前史节点。核心取证：《清穆宗实录》同治十一年相关条；户部奏销、盐政、河工或海关档案；同时期地方档案、地方志、日记或对方材料应按具体地区补检。该项锚定的是前史条件与信息背景，不把条件直接写成唯一因果。来源保留：官修编年和地方材料的记录密度与立场并不相同；本行不将条件、传闻、呈报或后见解释直接等同为确定因果。
\eventanchor{EQ-1872-CHINA-MERCHANTS-STEAM-NAVIGATION-DECISION}{围绕「轮船招商局创办，试图以…}同治十一年（1872年），围绕「轮船招商局创办，试图以官督商办方式发展近代航运」的命令、调度、照会或呈报形成独立的中央—地方行动文书链。核心取证：《清穆宗实录》同治十一年相关条；户部奏销、盐政、河工或海关档案。该项锚定的是命令或决策环节，命令抵达和结果仍须另外核验。来源保留：奏折、上谕、部议、军机档或条约可证明行政动作和机构立场，不自动证明所有地区、群体或对手按其意图行动。
\eventanchor{EQ-1872-CHINA-MERCHANTS-STEAM-NAVIGATION-AFTERMATH}{「轮船招商局创办，试图以官督…}同治十一年（1872年），「轮船招商局创办，试图以官督商办方式发展近代航运」引发的善后、执行或地方回应构成可由不同材料追踪的后续节点。核心取证：《清穆宗实录》同治十一年相关条；户部奏销、盐政、河工或海关档案；相关地区的档案、志书、契约、报刊或对方材料。该项锚定的是后续追踪问题，影响范围须按地区与群体分别判断。来源保留：战后、改革后或条约后的记忆、地方记录和官方总结常具有事后选择性；后续影响须按地点、群体和时间分层，不作整体化推断。
\eventanchor{EQ-1873-JAPAN-RYUKYU-DISPUTE-PRELUDE}{「日本围绕琉球与台湾事务加强…}同治十二年（1873年），「日本围绕琉球与台湾事务加强对清交涉」之前的条件、信息和争议已通过中央与地方材料显现，构成该事件可追踪的前史节点。核心取证：《清穆宗实录》同治十二年相关条；《筹办夷务始末》相关年卷与中外照会、条约文本；同时期地方档案、地方志、日记或对方材料应按具体地区补检。该项锚定的是前史条件与信息背景，不把条件直接写成唯一因果。来源保留：官修编年和地方材料的记录密度与立场并不相同；本行不将条件、传闻、呈报或后见解释直接等同为确定因果。
\eventanchor{EQ-1873-JAPAN-RYUKYU-DISPUTE-DECISION}{围绕「日本围绕琉球与台湾事务…}同治十二年（1873年），围绕「日本围绕琉球与台湾事务加强对清交涉」的命令、调度、照会或呈报形成独立的中央—地方行动文书链。核心取证：《清穆宗实录》同治十二年相关条；《筹办夷务始末》相关年卷与中外照会、条约文本。该项锚定的是命令或决策环节，命令抵达和结果仍须另外核验。来源保留：奏折、上谕、部议、军机档或条约可证明行政动作和机构立场，不自动证明所有地区、群体或对手按其意图行动。
\eventanchor{EQ-1873-JAPAN-RYUKYU-DISPUTE-AFTERMATH}{「日本围绕琉球与台湾事务加强…}同治十二年（1873年），「日本围绕琉球与台湾事务加强对清交涉」引发的善后、执行或地方回应构成可由不同材料追踪的后续节点。核心取证：《清穆宗实录》同治十二年相关条；《筹办夷务始末》相关年卷与中外照会、条约文本；相关地区的档案、志书、契约、报刊或对方材料。该项锚定的是后续追踪问题，影响范围须按地区与群体分别判断。来源保留：战后、改革后或条约后的记忆、地方记录和官方总结常具有事后选择性；后续影响须按地点、群体和时间分层，不作整体化推断。
\eventanchor{EQ-1874-JAPANESE-TAIWAN-EXPEDITION-PRELUDE}{「日本以牡丹社事件为由出兵台…}同治十三年（1874年），「日本以牡丹社事件为由出兵台湾南部，清廷被迫应对」之前的条件、信息和争议已通过中央与地方材料显现，构成该事件可追踪的前史节点。核心取证：《清穆宗实录》同治十三年相关条；军机处档、战事方略及地方奏报；同时期地方档案、地方志、日记或对方材料应按具体地区补检。该项锚定的是前史条件与信息背景，不把条件直接写成唯一因果。来源保留：官修编年和地方材料的记录密度与立场并不相同；本行不将条件、传闻、呈报或后见解释直接等同为确定因果。
\eventanchor{EQ-1874-JAPANESE-TAIWAN-EXPEDITION-DECISION}{围绕「日本以牡丹社事件为由出…}同治十三年（1874年），围绕「日本以牡丹社事件为由出兵台湾南部，清廷被迫应对」的命令、调度、照会或呈报形成独立的中央—地方行动文书链。核心取证：《清穆宗实录》同治十三年相关条；军机处档、战事方略及地方奏报。该项锚定的是命令或决策环节，命令抵达和结果仍须另外核验。来源保留：奏折、上谕、部议、军机档或条约可证明行政动作和机构立场，不自动证明所有地区、群体或对手按其意图行动。
\eventanchor{EQ-1874-JAPANESE-TAIWAN-EXPEDITION-AFTERMATH}{「日本以牡丹社事件为由出兵台…}同治十三年（1874年），「日本以牡丹社事件为由出兵台湾南部，清廷被迫应对」引发的善后、执行或地方回应构成可由不同材料追踪的后续节点。核心取证：《清穆宗实录》同治十三年相关条；军机处档、战事方略及地方奏报；相关地区的档案、志书、契约、报刊或对方材料。该项锚定的是后续追踪问题，影响范围须按地区与群体分别判断。来源保留：战后、改革后或条约后的记忆、地方记录和官方总结常具有事后选择性；后续影响须按地点、群体和时间分层，不作整体化推断。
\eventanchor{EQ-1874-TONGZHI-DEATH-PRELUDE}{「同治帝去世，慈禧选择光绪帝…}同治十三年（1874年），「同治帝去世，慈禧选择光绪帝继位，再次形成幼主政治」之前的条件、信息和争议已通过中央与地方材料显现，构成该事件可追踪的前史节点。核心取证：《清穆宗实录》同治十三年相关条；宫中朱批奏折、内阁大库题本与《清史稿》相关志传；同时期地方档案、地方志、日记或对方材料应按具体地区补检。该项锚定的是前史条件与信息背景，不把条件直接写成唯一因果。来源保留：官修编年和地方材料的记录密度与立场并不相同；本行不将条件、传闻、呈报或后见解释直接等同为确定因果。
\eventanchor{EQ-1874-TONGZHI-DEATH-DECISION}{围绕「同治帝去世，慈禧选择光…}同治十三年（1874年），围绕「同治帝去世，慈禧选择光绪帝继位，再次形成幼主政治」的命令、调度、照会或呈报形成独立的中央—地方行动文书链。核心取证：《清穆宗实录》同治十三年相关条；宫中朱批奏折、内阁大库题本与《清史稿》相关志传。该项锚定的是命令或决策环节，命令抵达和结果仍须另外核验。来源保留：奏折、上谕、部议、军机档或条约可证明行政动作和机构立场，不自动证明所有地区、群体或对手按其意图行动。
\eventanchor{EQ-1874-TONGZHI-DEATH-AFTERMATH}{「同治帝去世，慈禧选择光绪帝…}同治十三年（1874年），「同治帝去世，慈禧选择光绪帝继位，再次形成幼主政治」引发的善后、执行或地方回应构成可由不同材料追踪的后续节点。核心取证：《清穆宗实录》同治十三年相关条；宫中朱批奏折、内阁大库题本与《清史稿》相关志传；相关地区的档案、志书、契约、报刊或对方材料。该项锚定的是后续追踪问题，影响范围须按地区与群体分别判断。来源保留：战后、改革后或条约后的记忆、地方记录和官方总结常具有事后选择性；后续影响须按地点、群体和时间分层，不作整体化推断。
